\chapter{Deployment}

\begin{notebox}[NOTE]
This chapter covers the concepts and configuration values a production
MERN deployment needs, regardless of hosting provider. It does not cover
Git/GitHub workflow, provider-specific CLI push mechanics, or account
sign-up steps --- those vary by provider and are outside this handbook's
scope.
\end{notebox}

\section{Production Architecture}

\begin{diagrambox}[Production Request Path]
\begin{verbatim}
Internet
   |  HTTPS
   v
React Frontend (static hosting, e.g. CDN)
   |  HTTPS API calls (VITE_API_URL)
   v
Express Backend (Node hosting)
   |  MongoDB driver (TLS connection string)
   v
MongoDB Atlas (managed cluster)
\end{verbatim}
\end{diagrambox}

Two independent deployments, connected only by URLs and environment
variables --- the frontend is a static build artifact, the backend is a
long-running Node process, and the database is a separate managed service.

\section{Frontend: Build for Production}

\begin{lstlisting}
npm run build
# produces a dist/ folder: static HTML/CSS/JS, ready for any static host
\end{lstlisting}

The static host serves \texttt{dist/} directly --- no Node process runs
for the frontend in production. Set \texttt{VITE\_API\_URL} to the
deployed backend's URL \emph{before} running \texttt{build}, since Vite
inlines env variables at build time, not at runtime.

\begin{warnbox}[WARNING]
If \texttt{VITE\_API\_URL} still points at \texttt{localhost} when you
run \texttt{build}, every API call in production will try to reach the
visitor's own machine. Rebuild after changing the env value --- editing
\texttt{.env} alone does not update an already-built \texttt{dist/}.
\end{warnbox}

\section{Backend: Start Command}

\begin{lstlisting}
npm start
# package.json: "start": "node src/server.js"
\end{lstlisting}

Production runs \texttt{node}, never \texttt{nodemon} --- nodemon exists
purely for local file-watching and adds unnecessary overhead to a
production process. The hosting platform restarts the process on crash;
that is the platform's job, not nodemon's.

\section{Environment Variables}

Configure these on the hosting platform's dashboard/config system ---
never commit \texttt{.env} to the repository.

\begin{tabularx}{\textwidth}{@{}>{\bfseries}p{4.2cm}X@{}}
\toprule
Variable & Production value \\
\midrule
\texttt{MONGO\_URI} & Atlas connection string for the production cluster \\
\texttt{JWT\_SECRET} & Long random string, different from any dev secret \\
\texttt{NODE\_ENV} & \texttt{production} \\
\texttt{FRONTEND\_URL} & The deployed frontend's exact domain (for CORS) \\
\texttt{GOOGLE\_CLIENT\_ID/SECRET} & Production OAuth credentials with the
production redirect URI registered \\
\texttt{CLOUDINARY\_*} & Cloudinary account keys \\
\bottomrule
\end{tabularx}

\section{Updating CORS for Production}

\begin{lstlisting}[language=JavaScript]
app.use(cors({
  origin: process.env.FRONTEND_URL, // the deployed frontend's real domain
  credentials: true,
}));
\end{lstlisting}

A CORS origin of \texttt{localhost:5173} works locally but must be
replaced with the deployed frontend's actual domain, or every
cookie-based request will be blocked in production (Chapter 26).

\section{MongoDB Atlas Network Access}

Atlas rejects connections from IPs not on its allowlist. For a backend
hosted on a platform with dynamic/unknown outbound IPs, allow
\texttt{0.0.0.0/0} (any IP) --- acceptable because the connection still
requires valid database credentials. If the platform provides static
outbound IPs, allowlist those specific addresses instead for a tighter
boundary.

\section{HTTPS and Custom Domains}

Most hosting platforms provision HTTPS automatically (managed
certificates) for both the frontend and backend as soon as a domain is
attached --- no manual certificate handling is typically required. Once
HTTPS is active, set cookie options to \texttt{secure: true} and
\texttt{sameSite: 'none'} for cross-domain cookie delivery (Chapter 19).
A custom domain is simply a DNS record pointed at the platform, applied
to either the frontend, backend, or both independently.

\section{Debugging in Production}

Every hosting platform exposes a logs view for the running backend
process --- this is the production equivalent of the terminal output you
see locally. When something fails only in production:

\begin{itemize}
\item Reproduce the request and immediately check the platform's log
viewer for the stack trace.
\item Confirm the failing behavior isn't simply a missing/incorrect
environment variable (the most common production-only bug).
\item Check the browser console and Network tab on the deployed frontend
exactly as you would locally.
\end{itemize}

\whatwhyhow
{Two independently deployed pieces --- a static frontend build and a
long-running Node backend process --- connected through environment
variables (\texttt{VITE\_API\_URL}, \texttt{FRONTEND\_URL},
\texttt{MONGO\_URI}) and a managed MongoDB Atlas cluster.}
{Separating concerns lets the frontend scale as static assets on a CDN
while the backend scales as a process, and keeps secrets out of the
repository entirely.}
{Build the frontend with the production API URL baked in, start the
backend with \texttt{npm start} (not nodemon), configure CORS/cookies for
the real domains, and set every secret through the platform's environment
variable system.}
